\section{The Shattering of the Vessels and the Sparks}

The Lurianic creation account carries a \textbf{primordial break} built into the structure of the
cosmos. Before the present ordered world there was an earlier, unstable world that could not hold the
light. It shattered. From that shattering came the shells of evil and the holy sparks trapped inside
them. This section gives that catastrophe and what it left behind.

\subsection{Olam ha-Tohu, the World of Chaos}

The \textbf{World of Chaos} (Olam ha-Tohu), named from the tohu va-vohu of Genesis, was the
first emanated world.

\begin{itemize}
 \item It had \textbf{abundant light} but \textbf{weak vessels}.
 \item The ten lights stood apart in a straight line. Each acted alone with no give-and-take. Each
 tried to receive the full light by itself.
 \item Strong light plus weak, unconnected vessels is an unstable combination. It could not hold.
\end{itemize}

This is contrasted with the \textbf{World of Repair} (Olam ha-Tikun), where the lower lights sit in
strong, interconnected vessels and so endure.

\subsection{Shevirat ha-Kelim, the Shattering}

The weak vessels of Tohu could not contain the strong light. They \textbf{shattered}. This is the
\textbf{breaking of the vessels} (Shevirat ha-Kelim).

The seven lower vessels broke apart, along with parts of the upper ones. The break is the origin of
disorder, exile, and death in the cosmos. It is a catastrophe built into creation itself.

\subsection{The Kings of Edom}

The shattering is also called the \textbf{death of the kings of Edom}. Genesis 36 lists eight kings of
Edom who reigned before any king reigned over Israel, and of several it says simply that each one died.
The Ari reads each death as the breaking of a vessel and the falling of its sparks.

\begin{itemize}
 \item The \textbf{seven kings who died} are the seven shattered vessels of Tohu, the seven emotional
 sefirot standing alone.
 \item The \textbf{eighth king, Hadar}, is described differently. His death is not stated, and the name
 of his wife is given. This signals the beginning of repair, the pairing of male and female. Hadar is
 the turning point from Tohu to Tikun.
\end{itemize}

\subsection{The Qelippot, the Shells}

When the vessels shattered, the broken shards fell and became the \textbf{shells} (qelippot), the
husks. They form the \textbf{Other Side} (Sitra Achra), the impure counter-structure of evil.

The image is a fruit. The edible fruit is the holiness, the divine light. The shell is the inedible peel
that surrounds and conceals it. The shells do three things.

\begin{itemize}
 \item They \textbf{conceal} the divine light. Evil is, at root, hiddenness.
 \item They \textbf{trap} sparks of holiness inside themselves.
 \item They \textbf{feed} on the light they imprison. The Other Side draws its life from the holiness
 it conceals.
\end{itemize}

Among the shells one is special. The \textbf{translucent shell} (Qelippat Nogah) is neutral and mixed,
partly impure and partly able to be elevated. It is the realm of free choice, where the trapped sparks
can be redeemed. The other three shells are wholly impure and can only be rejected.

\subsection{The Sparks (Netzotzot)}

Clinging to each fallen shard is a \textbf{spark of holy light}, a \textbf{nitzotz} (plural
\textbf{netzotzot}). Holy sparks are now trapped inside the husks, exiled, scattered through the lower
worlds and through matter.

The light is still holy. It is simply imprisoned in the wrong place. So the central task of creation is
to \textbf{sort} the spark from the shell and \textbf{raise} it back to its source. Every time a person
uses a physical thing for a holy purpose, they release the spark trapped in it.

\subsection{Evil as Broken Structure}

This is the key conceptual point. Evil is \textbf{not a thing God made directly}. It is broken structure
plus trapped light. Evil is:

\begin{itemize}
 \item \textbf{Derivative, not original.} It arises as a consequence of the structure, not from a rival
 source. There is one root, Ein Sof.
 \item \textbf{Parasitic, not self-sustaining.} The Other Side lives only off the sparks it traps. Cut
 off the stolen light and it cannot exist.
 \item \textbf{Concealment, not substance.} At bottom evil is hiddenness of the good, a husk over the
 light.
 \item \textbf{Doomed, not eternal.} When all the sparks are gathered and the shells are emptied, the
 Other Side simply collapses.
\end{itemize}

This separates Kabbalah from \textbf{strict dualism}. Strict dualism, as in some Gnostic or
Zoroastrian systems, posits two independent first principles, a good god and an evil god, locked in
eternal war as equals. Kabbalah rejects this. The imagery shows two opposed sides, the side of holiness
and the Other Side, but they are not metaphysical equals. One is the source. The other is a temporary,
dependent, doomed shadow that exists to be overcome. Darkness is not defeated in battle. It is drained
of its borrowed light. The Zoroastrian contrast is taken up in its own chapter.
